\chapter{Patterns: feature flags}
\label{ch:flags}

\begin{aside}
Ship code that's behind a flag. Enable for beta users. Roll
out gradually. Roll back instantly. The standard pattern in
modern software.
\end{aside}

\section{Flag definition}

\begin{lstlisting}
struct Flags {
    Signal<bool> dark_mode = false;
    Signal<bool> new_search = false;
    Signal<bool> experimental_charts = false;
};
\end{lstlisting}

Flags as signals enable runtime toggle.

\section{Usage}

\begin{lstlisting}
auto search = ctx.get(flags.new_search)
            ? new_search_widget()
            : old_search_widget();
\end{lstlisting}

\section{Sources}

Flags come from:

\begin{itemize}[leftmargin=*]
  \item Config file.
  \item Environment variables.
  \item Command-line arguments.
  \item Remote config service.
\end{itemize}

\section{Cohort selection}

For gradual rollout: a flag is on if your user hash falls
in the active range. \code{hash(user\_id) \% 100 < N}.

\section{Cleanup}

A feature graduates: remove the flag, keep only the new
code. Tech-debt budget for this.

\section{Recap}

\begin{recap}
\begin{itemize}[leftmargin=*]
  \item Flags as signals for runtime toggle.
  \item Multiple sources (file, env, remote).
  \item Cohort selection for gradual rollout.
  \item Clean up old flags.
\end{itemize}
\end{recap}

\section*{Workshop}

\begin{enumerate}[leftmargin=*]
  \item Add a feature flag for one experimental
    feature.
  \item Wire to a config file.
  \item Plan its cleanup.
\end{enumerate}
